\section{本章小结}
\label{sec:wave-summary}

本章是椭圆方程和扩散方程之后的补充，讨论二阶时间演化的基本认识。由弦的受力得到波动方程，弯曲决定加速度，所以初始状态是位移和速度一对函数。

一维实轴上，d'Alembert公式把解写成一列右行波和一列左行波。初始位移分成两半走开，初始速度留下不断变宽的平台，依赖区域由两条特征线截出。端点通过反射起作用：固定端翻转，自由端不翻转。

区间上的解按边界条件允许的模式展开，时间因子是不衰减的简谐振动，驻波正是两列行波的叠加。外力用Duhamel原理化成脉冲的叠加，模式频率与外力匹配时出现共振，边界驱动也按功率进入能量账。

表~\ref{tab:wave-summary}对比了扩散方程和波动方程。

\begin{BookInlineTable}
  \centering
  \caption{扩散方程与波动方程的对比}
  \label{tab:wave-summary}
  \begin{tabular}{lll}
    \toprule
    & 扩散方程 & 波动方程\\
    \midrule
    时间阶数 & 一阶 & 二阶\\
    初值 & $u_0$ & $u_0,u_1$\\
    模式因子 & 指数衰减 & 简谐振荡\\
    能量 & $\frac12\lVert u\rVert_2^{2}$单调减小 & 动能加弹性能守恒\\
    传播速度 & 无穷（尾部极弱） & 不超过$c$\\
    折角 & 立即消失 & 沿特征线传播\\
    时间反向 & 不适定 & 适定\\
    外力下 & 无共振 & 共振时线性增长\\
    \bottomrule
  \end{tabular}
\end{BookInlineTable}

能量恒等式不依赖解的公式，由它得到守恒、唯一性和有限时间的增长估计；局部能量进一步给出有限传播速度，在任意维成立。能量守恒不等于振幅不变，前沿之外为零也不说明前沿之内的形状相同，这两点要与一维公式分开理解。

本章只讨论常系数的标量波动方程和有限时间的结论。变系数介质、弹性力学和电磁学中的波动系统，要在具体物理建模需要时，再按一般的双曲系统扩展。
